\documentclass[11pt]{article}
\usepackage[margin=2.7cm]{geometry}
\usepackage[T1]{fontenc}
\usepackage[utf8]{inputenc}
\usepackage{lmodern,amsmath,amssymb,amsthm}
\usepackage[hidelinks]{hyperref}
\newtheorem{proposition}{Proposition}
\theoremstyle{remark}\newtheorem*{remark}{Remark}
\newcommand{\Z}{\mathbb{Z}}
\title{One reversible measure, but a non-reversible\\ stationary measure: exclusion processes}
\author{counterexamples-ips, entry 0001}
\date{Status: \emph{plausible} (not yet checked by a human)\\[2pt] \small Idea and write-up: Claude (Opus 5.5), 2 November 2026}
\begin{document}
\maketitle

\paragraph{Question.} Let an interacting particle system on $\{0,1\}^{\Z^d}$ have translation-invariant, finite-range rates. If it has at least one reversible measure, is every stationary measure reversible?

For stochastic Ising models this is the Holley--Stroock question: true for translation-invariant stationary measures in all $d$ (Holley 1971), true for all stationary measures when $d\le 2$ (Holley--Stroock 1977), open for $d\ge3$. In the generality above, the answer is \emph{no}.

\paragraph{Setting.} Nearest-neighbour asymmetric simple exclusion (ASEP) on $\Z$: particles jump right at rate $p$ and left at rate $q$, $p>q>0$, jumps to occupied sites being suppressed:
\[
\Omega f(\eta)=\sum_{x\in\Z}\Big(p\,\eta(x)\big(1-\eta(x+1)\big)+q\,\eta(x+1)\big(1-\eta(x)\big)\Big)\big(f(\eta^{x,x+1})-f(\eta)\big).
\]
Let $\lambda=p/q>1$.

\begin{proposition}
For ASEP on $\Z$ with $p>q$:
\begin{enumerate}
\item for every $c>0$, the product measure $\mu_c$ with $\mu_c(\eta(x)=1)=\rho_x:=c\lambda^x/(1+c\lambda^x)$ (a \emph{blocking measure}) is reversible;
\item for every $\rho\in(0,1)$, the Bernoulli product measure $\nu_\rho$ is stationary but not reversible.
\end{enumerate}
\end{proposition}

\begin{proof}
(1) For product measures and exclusion dynamics it suffices to check detailed balance on each bond (Liggett 1985, Ch.~VIII): $p\,\rho_x(1-\rho_{x+1})=q\,\rho_{x+1}(1-\rho_x)$. Since $\rho_y/(1-\rho_y)=c\lambda^y$, the ratio $\rho_{x+1}(1-\rho_x)/\bigl(\rho_x(1-\rho_{x+1})\bigr)$ equals $\lambda=p/q$.

(2) Stationarity of $\nu_\rho$ holds for every translation-invariant exclusion process (Liggett 1985, Ch.~VIII). For non-reversibility, let $f(\eta)=\eta(0)(1-\eta(1))$ and $g(\eta)=(1-\eta(0))\eta(1)$. On $\{f=1\}$ we have $g=0$, and the only transition producing $g=1$ is the jump $0\to1$, at rate $p$; all other transitions touching $0$ or $1$ are blocked or lead to $(\eta(0),\eta(1))\in\{(0,0),(1,1)\}$. Hence $f\,\Omega g=pf$, and likewise $g\,\Omega f=qg$, so
\[
\int f\,\Omega g\,d\nu_\rho=p\rho(1-\rho)\neq q\rho(1-\rho)=\int g\,\Omega f\,d\nu_\rho .\qedhere
\]
\end{proof}

\paragraph{Mechanism.} Reversibility means zero net current across every bond. Under $\nu_\rho$ the current is $(p-q)\rho(1-\rho)>0$: a translation-invariant state with drift transports particles. The blocking measures pile the particles up ($\rho_x\to0$ as $x\to-\infty$, $\rho_x\to1$ as $x\to+\infty$), so that exclusion exactly balances the drift on every bond. Conservation of particles lets both kinds of stationary state coexist in sectors of the state space that the dynamics never connects.

\begin{remark}[Degenerate version]
For any exclusion process $\delta_{\mathbf 0}$ and $\delta_{\mathbf 1}$ are reversible, so the claim fails for asymmetric exclusion in every dimension. The blocking measures show the failure is not merely due to absorbing configurations.
\end{remark}

\begin{remark}[Full support, $d\ge2$; sketch]
On $\Z^2$ let particles jump to $x\pm e_1$ at rates $p,q$ ($p>q$) and to $x\pm e_2$ at rate $r$ each. The product measure with $\rho_x/(1-\rho_x)=\lambda^{x_1}$ satisfies detailed balance on every bond, so it is reversible with full support, while $\nu_\rho$ is stationary and not reversible.
\end{remark}

\paragraph{Which hypothesis breaks.} The claim must exclude conservative dynamics. In the Holley--Stroock setting the rates are positive spin-flip rates, so there is no conservation law and no decomposition into sectors.

\paragraph{Follow-up questions.} See \texttt{entries/0001-asep-blocking-vs-product-measures/followups.md}.

\begin{thebibliography}{9}\small
\bibitem{H71} R.~Holley, Free energy in a Markovian model of a lattice spin system, \emph{Comm.\ Math.\ Phys.}\ 23 (1971).
\bibitem{HS77} R.~Holley, D.~Stroock, In one and two dimensions, every stationary measure for a stochastic Ising model is a Gibbs state, \emph{Comm.\ Math.\ Phys.}\ 55 (1977).
\bibitem{L76} T.~M.~Liggett, Coupling the simple exclusion process, \emph{Ann.\ Probab.}\ 4 (1976).
\bibitem{L85} T.~M.~Liggett, \emph{Interacting Particle Systems}, Springer, 1985, Ch.~IV and VIII.
\bibitem{L99} T.~M.~Liggett, \emph{Stochastic Interacting Systems: Contact, Voter and Exclusion Processes}, Springer, 1999, Part~III.
\end{thebibliography}
\end{document}
